\documentclass[11pt]{article}
\usepackage[margin=0.85in]{geometry}
\usepackage{amsmath,amssymb,booktabs,array,longtable}
\usepackage[T1]{fontenc}
\usepackage{lmodern,microtype,xcolor}
\usepackage[hidelinks]{hyperref}
\usepackage{fancyhdr}
\definecolor{navy}{HTML}{173957}
\pagestyle{fancy}\fancyhf{}\lhead{Gamma-fitness tMN study}\rhead{Four-year protocol}\cfoot{\thepage}
\setlength{\headheight}{14pt}
\setlength{\parindent}{0pt}\setlength{\parskip}{6pt}
\setlength{\emergencystretch}{2em}
\newcommand{\tmn}{\mathrm{tMN}}
\title{\color{navy}A four-year study of clonal evolution\texorpdfstring{\\}{ }and tMN after CAR-T therapy}
\author{Study protocol and implementation specification}
\date{5 November 2026}
\begin{document}
\maketitle
\textbf{Purpose.} Replace the previous LD-plus-exponential fitness model with a gamma-shaped CAR-associated pulse, fit only its shape and scale, jointly use longitudinal VAF and clinical tMN information, repeat the fit under fixed-parameter sensitivity scenarios, and estimate cumulative model risks at 1, 2, 3 and 4 years after CAR infusion.

\textbf{Implementation decision.} The instruction to begin is taken to authorize the proposed reference configuration: peak-normalized gamma fitness, one shared $(\alpha,\theta)$ pair, a patient-level index-marker analysis, and blood-cell longitudinal observations where available. Reference values below are modeling assumptions, not measured biological constants. The alternative unit-area gamma specification is documented rather than silently treated as equivalent.

\section{Population process and time origin}
CAR infusion is day 0. One year is 365 days and the analysis horizon is day 1460. Retain the fixed-size Moran process with
\[
W(t)+M(t)+D(t)=N.
\]
$W$ lacks the modeled founder marker; $M$ carries that marker but has not acquired a second transforming alteration; $D$ aggregates cells acquiring that alteration and their descendants. Multiple independent second-hit origins are aggregated. This is an operational two-hit model, not a claim that two specific molecular lesions universally cause tMN.

Replacement events occur at rate $Nb$. With $w=N-m-d$, $F=w+f_Mm+f_Dd$, newborn probabilities are
\[
p_W=\frac wF,\qquad p_M=\frac{(1-\mu)f_Mm}{F},\qquad
p_D=\frac{f_Dd+\mu f_Mm}{F}.
\]
The off-diagonal generator is $q_{n,n+e_j-e_k}=b n_kp_j$ for $j\ne k$. Here $\mu$ is the second-hit probability per $M$ division, not a daily probability.

\textbf{Remove LD fitness completely.} For $t<0$, set
\[
f_W=f_M=f_D=1.
\]
Neutral turnover, drift and mutation may still occur. No LD fitness tail survives after day 0. The fixed-$N$ framework does not introduce a separate LD population bottleneck.

\section{Gamma-associated fitness after CAR infusion}
Use the shape--scale gamma density
\[
g(t;\alpha,\theta)=\frac{t^{\alpha-1}e^{-t/\theta}}
{\Gamma(\alpha)\theta^\alpha},\quad t>0,\quad\alpha>1,\quad\theta>0.
\]
$\alpha$ is dimensionless; $\theta$ is a scale in days, not a rate. Its mode is $t_p=(\alpha-1)\theta$. The new $\alpha$ is unrelated to the previous exponential-rise rate with the same symbol.

\textbf{Primary specification: fixed peak.}
\[
h(t)=\begin{cases}0,&t\le0,\\g(t)/g(t_p),&t>0,\end{cases}
\qquad
f_M(t)=1+A_{\rm peak}h(t),\quad
f_D(t)=f_M(t)+s_D\quad(t\ge0).
\]
Only $\alpha,\theta$ are fitted. $A_{\rm peak}$ is fixed and dimensionless. The model starts with zero CAR-associated advantage and permits a rise and fall. The extra $D$ advantage starts at day 0; it is zero before infusion.

\textbf{Alternative: fixed area.} A literal normalized gamma density requires
$f_M=1+A_{\rm area}g(t)$, where $A_{\rm area}$ has units of days and fixes the integrated advantage. Directly adding the PDF in day units implicitly sets $A_{\rm area}=1$ day. In the two-type, no-mutation limit, the cumulative log-odds increment is bounded above by $bA_{\rm area}$; at $b=1/7$ and $A_{\rm area}=1$, the odds can increase at most by $e^{1/7}\simeq1.154$. Fixed area and fixed peak are different biological restrictions. See the \href{https://www.itl.nist.gov/div898/handbook/eda/section3/eda366b.htm}{NIST gamma distribution reference} for the density parameterization.

\section{Reference parameters and sensitivity scenarios}
\begin{longtable}{@{}p{0.13\textwidth}p{0.39\textwidth}p{0.17\textwidth}p{0.23\textwidth}@{}}
\toprule
Parameter & Definition & Reference & Sensitivity values\\\midrule\endhead
$N$ & Effective competing population size; not total body hematopoietic cells & $10^5$ (retained) & $10^4,10^5,10^6$\\
$b$ & Moran turnover scale, day$^{-1}$ & $1/7$ (retained) & $1/30,1/7,1$\\
$A_{\rm peak}$ & Maximum extra founder fitness & $0.1$ (proposed) & $0.03,0.1,0.3$\\
$\mu$ & Second hit per $M$ division & $10^{-6}$ (proposed) & $0,10^{-7},10^{-6},10^{-5},10^{-4}$\\
$r_D$ & Additional rare-clone growth scale $bs_D$, day$^{-1}$ & $0.003$ (proposed) & $0,0.001,0.003,0.01$\\
$s_D$ & Extra post-CAR fitness of $D$ over $M$ & $r_D/b=0.021$ & Derived; hold $r_D$ fixed when varying $b$\\
$D_{\rm init}$ & Double-hit cells at the initial observed baseline & 0 (proposed) & 1, 10, 100; compatible with founder burden\\
$\kappa$ & Operational threshold fraction & $0.01$ (retained) & $0.005,0.01,0.02$\\
$\sigma_{\rm VAF}$ & Logit measurement SD if read counts unavailable & 0.25 (proposed) & $0.1,0.25,0.5$\\
Delay & Threshold-to-diagnosis delay, days & 0 (simplification) & Evidence-based extensions\\
\bottomrule
\end{longtable}
The old study specified grids, not a unique reference, for $\mu,r_D,D_{\rm init}$. The new reference values select explicit assumptions from those grids. $r_D$ is an approximate rare-clone growth scale, not the exact net rate in a mixed population. Set $K=\lfloor\kappa N\rfloor+1$ whenever $N$ or $\kappa$ changes. $K=1001$ in the reference. Remove LD parameters and independent exponential rise/decay parameters.

Numerical search bounds are initially $1.01\le\alpha\le50$, $0.25\le\theta\le1460$ days. A boundary optimum triggers review and is not evidence that a parameter is identified. Derived quantities include peak day, pulse duration and integrated advantage.

\section{Cohort construction and longitudinal VAF}
Retain the six-gene scope from the latest study: DNMT3A, PPM1D, TP53, TET2, CBL and ASXL1. The supplied files are MOESM7 and MOESM9 of article identifier 41588\_2026\_2795. Preserve source hashes, worksheet names and Excel row numbers. Match exact patient, locus and allele pair; record orientation checks and unresolved or duplicate matches.

Use all available prediagnosis longitudinal measurements within the four-year window, rather than only first/last VAF. Do not pool blood cells, cfDNA and tumor samples. Prefer PBL or PBMC observations, separately by sample type. Sample-level depth is not automatically a locus-specific binomial denominator.

\textbf{Index-marker approximation.} In the absence of phasing, select one marker per patient using the largest eligible baseline VAF at or before day 0, with deterministic gene/variant tie-breaking. Do not select by later expansion or tMN outcome. Track this same variant and sample type. Record all other variants as secondary data; do not sum their VAF or risks. This assumes the selected founder can represent the modeled transformation pathway and does not establish that it caused the clinical tMN. A model using every unphased variant needs an explicit multiclone architecture.

Require a positive, diploid-compatible baseline and a later eligible observation for the initial VAF diagnostic fit. Patients with only post-infusion entry need a left-entry model, not a substituted day-0 baseline. Condition on the measured baseline for initialization; this initial implementation does not propagate its measurement error. Set total founder burden to $2Nv_{i0}$ and subtract $D_{\rm init}$ to obtain $M_{\rm init}$. Initial $D$ must be below $K$ and no greater than founder burden. Baseline can be at a negative day and the process is propagated neutrally to day 0.

Under a diploid, heterozygous, representative-sample assumption,
\[
v(t)=\frac{M(t)+D(t)}{2N}.
\]
If valid alternate and total read counts exist, use a count observation model. Otherwise use
\[
z_{ij}=\operatorname{logit}(2v_{ij}),\quad
z_{ij}\mid X_i(t_{ij})\sim
N\!\left(\operatorname{logit}\frac{M_i+D_i}{N},\sigma_{\rm VAF}^2\right).
\]
The variance is fixed, not fitted. Zero/non-detected values require an explicit detection model; do not invent read counts or pseudocounts. Flag $v\ge0.5$, LOH/CNV concerns and incompatible compartments. Excluding non-detections makes the initial analysis conditional on positive detections; this selection limitation must accompany results. Exclude observations at or after tMN diagnosis from the prediagnosis process. Unknown treatment changes remain a limitation.

Patients with unresolved multiple CAR origins or source-positive tMN without a diagnosis time are withheld from the initial primary VAF diagnostic. Sample-day metadata conflicts are excluded; sample-type naming conflicts remain visible for review.

\section{Clinical outcomes and data readiness}
Each patient contributes one clinical outcome. Establish the index infusion, first diagnosis, last verified event-free contact and death. Multiple CAR origins must be reconciled. Last VAF sampling is not a clinical censor date. Absence from the source tMN sheet is not a verified negative.

Current extracted evidence has no verified event-free follow-up dates. Source-positive diagnoses include events beyond four years; some times or origins are missing or inconsistent. Build a reconciliation table before a clinical fit. A later event supports censoring at day 1460 only after the first-event status and time origin are verified. Keep unresolved labels missing. Do not use case-only outcomes from an enriched ascertainment scheme as an unqualified population-incidence calibration.

\section{Joint likelihood}
Let $\psi$ contain every fixed parameter. With $Y_i$ the VAF series and $O_i$ the event/censor record, use
\[
L(\alpha,\theta\mid\psi)=\prod_i\int
p(X_i\mid\alpha,\theta,\psi)\,
p(Y_i\mid X_i,\psi)\,
p(O_i\mid X_i,\psi)\,dX_i.
\]
VAF and tMN depend on the same latent process. Multiplying separately computed marginal VAF and marginal event probabilities would discard that dependence. An exact event time uses boundary flux; right censoring uses survival; interval censoring uses a survival difference; a verified status at a specified horizon uses its Bernoulli probability. Missing outcomes contribute no clinical outcome factor.

Fit one cohort-wide $(\alpha,\theta)$ pair initially. Patient--gene-specific two-parameter fits from a single endpoint change are generally weakly identified. Gene-specific shared parameters are an extension requiring sufficient longitudinal support, not the primary model. Inspect profile likelihoods and synthetic parameter recovery before interpreting fitted values.

\textbf{Initial numerical implementation.} Let $q=M/(N-D)$ be a deterministic background and evolve
\[
\dot q=bq(1-q)\frac{(1-\mu)f_M-1}{1+q(f_M-1)}.
\]
Given $q$, evolve a stochastic $D$ birth--death chain with $m=(N-d)q$ and rates
\[
\lambda_d=b(N-d)p_D,\qquad \delta_d=bd(1-p_D).
\]
Apply VAF observation weights directly to this distribution, normalize and accumulate log normalizers. For known pre-event/censored records use the killed chain $d<K$; an event contributes its boundary flux. For unknown outcomes use the unrestricted physical chain $d=0,\ldots,N$, rather than implicitly conditioning on no crossing. This hybrid approximates fluctuations of $M$ and their feedback and requires validation. The reference threshold proxy is not automatically equivalent to clinical diagnosis.

A separate coupled mean-field VAF-only fit can initialize the search and expose incompatibilities while outcome reconciliation is pending. It retains $\mu,D_{\rm init}$ and $s_D$, but is explicitly a preliminary diagnostic, not the completed joint likelihood. No final clinical risk fit is substituted for missing follow-up.

The initial forward-likelihood implementation uses positive backward-Euler steps and must be time-step refined. For standalone risk curves, an adaptive banded LSODA solver with an explicit absorbing-flux state reduces relative error in very small early-event probabilities. Validate tolerance refinements separately from the deterministic-background approximation.

\section{Sensitivity, uncertainty and four-year risks}
For each fixed-parameter scenario: reconstruct compatible initial states, refit $\alpha,\theta$ using the same joint likelihood, check fit and identifiability, and recalculate risks. Begin with one-at-a-time scenarios; selected interactions include $\mu\times r_D$, $\mu\times D_{\rm init}$ and $A_{\rm peak}\times b$. Keep the chosen marker/cohort fixed for comparable scenarios and flag any incompatibility rather than silently dropping patients.

Define $T_*=\inf\{t\ge0:D(t)\ge K\}$. For initially event-free patients, report
\[
P(T_*\le365),\quad P(T_*\le730),\quad
P(T_*\le1095),\quad P(T_*\le1460).
\]
If baseline precedes infusion, condition the day-0 distribution on not having crossed before infusion. Report patient-level risks and clearly defined cohort summaries. Do not add gene risks. Separate patient-bootstrap/likelihood uncertainty, fixed-parameter scenario ranges and numerical error. A $\mu=D_{\rm init}=0$ scenario gives zero event risk and is incompatible with a positive event; never hide this using a probability floor.

The 1\% threshold is an operational proxy. Without a clinical observation mapping and competing mortality, label results as model first-passage risks with death omitted. Real-world clinical cumulative incidence requires death as a competing event. A retrospective fit using later VAF is not an independently validated prospective day-0 prediction. Source-cohort selection also limits generalization.

\section{Validation, execution and deliverables}
\begin{enumerate}
\item Audit longitudinal joins, allele orientation, duplicates, sample types, baseline timing, VAF range, diagnosis and censoring. Preserve unresolved records and supply a reconciliation template.
\item Verify pre-CAR neutrality, gamma peak/area behavior, zero-mutation invariants, probability conservation and monotone cumulative risks. Compare time-step refinements and a small-$N$ exact process; perform synthetic joint-likelihood recovery.
\item Run an explicitly labeled preliminary VAF diagnostic fit while auditing outcomes. Inspect residuals, selection boundaries and $(\alpha,\theta)$ likelihood ridges.
\item Once outcome readiness is established, run the clinical joint fit, patient-level uncertainty analysis and complete refitted sensitivity scenarios. If readiness fails, report the exact missing fields and withhold an unsupported clinical probability claim.
\item Deliver English protocol TeX/PDF, versioned configuration, source/data audit, reusable scripts, fit diagnostics, sensitivity tables, and 1--4-year CSV/XLSX risk tables and figures. Keep this study separate from previous versions.
\end{enumerate}
\textbf{Reproducibility.} Record package versions, input hashes, numerical settings, seeds, execution status and timestamps. A successful synthetic check or preliminary VAF fit must not be reported as completion of clinical joint inference.
\IfFileExists{initial_status.tex}{\begingroup\small\clearpage
\appendix
\section{Initial execution record: 5 November 2026}
\textbf{Status.} Protocol authoring, data extraction, initial engine validation and preliminary VAF diagnostics have been executed. Clinical joint inference is pending verified outcome and follow-up data. No missing clinical outcome was coded as a negative.

\textbf{Source and cohort audit.} Scanned 135,806 mutation rows and matched 1,696 longitudinal records to the selected six-gene endpoint inventory: 973 cfDNA, 353 PBL, 320 PBMC and 50 tumor observations. All matched endpoint allele pairs use REF-to-TUMOR orientation. The audit flags 98 duplicate-call rows, 287 values outside the positive diploid mapping, and 320 sample-type naming conflicts. These flag categories can overlap. There are no sample-day metadata conflicts among these matched observations.

The initial patient-level blood-cell diagnostic includes 12 patients with 54 observations: 12 baseline initial conditions and 42 likelihood observations. Exclusions among 67 source patients are 21 without a unique eligible pre-CAR blood-cell baseline, 31 whose baseline-selected marker has no eligible follow-up, 2 with unresolved multiple CAR time origins, and 1 source-positive patient without a diagnosis time. These stricter rules materially narrow the earlier endpoint cohort. The 320 raw PBMC records are labeled PBL in sample metadata; the mutation-table label is retained and compartments are not pooled within a patient.

\textbf{Validation.} Initial invariants passed. An independent constant-rate matrix exponential agrees with the refined shared-state joint log likelihood to 0.0068 log units at a 0.125-day backward-Euler step. One small-$N$ full Moran comparison differs from the hybrid risk by approximately $1.11\times10^{-4}$ in absolute probability. Noiseless synthetic mean-field data recover known shape/scale parameters. This is not a clinical joint parameter-recovery experiment. Adaptive full-$N$ risk tolerance checks also passed. These are initial numerical checks, not broad certification of the hybrid approximation.

\textbf{Preliminary fit.} One reference and eight individually refitted VAF-only scenarios have completed. The reference has
$\alpha=43.560$, $\theta=0.3708$ days, and peak day $15.78$. The standardized residual sum of squares is $624.78$ across 42 observations, with maximum absolute residual $10.86$ measurement SDs. The shared nonnegative pulse leaves substantial lack of fit; some sensitivity fits also reach the numerical search boundary. Optimizer convergence does not establish biological adequacy or parameter certainty. No clinical outcomes were used and no confidence intervals have been established.

\textbf{Risk calculation demonstration only.} The following are model threshold probabilities under the preliminary VAF fit and fixed reference parameters, starting at day 0 with the indicated VAF. They are \textbf{not clinically calibrated tMN probabilities}, omit competing death, and are not patient-specific estimates. Entries are probabilities on a 0--1 scale, not percentages.

\begin{center}
\begin{tabular}{@{}lrrrr@{}}
\toprule
Baseline VAF & Year 1 & Year 2 & Year 3 & Year 4\\\midrule
0.5\% & $4.34\times10^{-9}$ & $2.67\times10^{-5}$ & $3.52\times10^{-4}$ & $0.00111$ \\
1\% & $8.7\times10^{-9}$ & $5.35\times10^{-5}$ & $7.03\times10^{-4}$ & $0.00223$ \\
5\% & $4.41\times10^{-8}$ & $2.69\times10^{-4}$ & $0.00351$ & $0.011$ \\
\bottomrule
\end{tabular}
\end{center}
Standalone risks use adaptive banded LSODA with explicit absorbed flux and a tenfold tolerance comparison. The maximum absolute difference across these demonstrations is below $3\times10^{-10}$. This numerical agreement does not resolve the clinical data and model assumptions.

\textbf{Remaining clinical requirement.} Verify index infusion, first diagnosis or event-free censor date, death information and ascertainment. The readiness check currently finds no analysis patient with a complete verified clinical record. The shared-state likelihood implementation is available, but the clinical optimizer, full refitted clinical sensitivity study, uncertainty intervals and final clinical 1--4-year estimates have not been run. See the CSV reconciliation table and machine-readable readiness/status reports.
\endgroup}{}
\end{document}
